\section{Step 1: Basic Project Skeleton}

Every non-trivial program starts as something almost trivial. Before sockets,
TLS, or argument parsing, we need a file that compiles, accepts a URL on the
command line, and proves the tool is wired up.

\subsection{Why This Step Exists}

If \texttt{main} cannot see \texttt{argv}, nothing later matters. A skeleton
also gives you a compile command you will reuse for the rest of the tutorial:

\begin{lstlisting}[style=snippet,language=bash]
gcc -Wall -Wextra -O2 -o urlget urlget.c
./urlget http://example.com
\end{lstlisting}

\subsection{What We Add}

\begin{itemize}[leftmargin=1.4em]
    \item Standard headers for I/O and \texttt{EXIT\_FAILURE}
    \item A check that at least one argument (the URL) was provided
    \item A confirmation print so you can see the argument was received
\end{itemize}

\subsection{New Code}

\begin{lstlisting}[style=snippet,language=C]
#include <stdio.h>
#include <stdlib.h>

int main(int argc, char *argv[])
{
    if (argc < 2) {
        fprintf(stderr, "Usage: %s <URL>\n", argv[0]);
        return EXIT_FAILURE;
    }
    printf("Requested: %s\n", argv[1]);
    return 0;
}
\end{lstlisting}

Error messages go to \texttt{stderr}; successful chatter goes to
\texttt{stdout}. That split stays useful once the program starts printing
HTTP bodies to stdout.

\FullSourceStart{urlgetstep01}
\begin{lstlisting}[style=newcode,language=C,firstnumber=1]
/* urlget.c -- Step 1: basic skeleton */
#include <stdio.h>
#include <stdlib.h>

int main(int argc, char *argv[])
{
    if (argc < 2) {
        fprintf(stderr, "Usage: %s <URL>\n", argv[0]);
        return EXIT_FAILURE;
    }
    printf("Requested: %s\n", argv[1]);
    return 0;
}
\end{lstlisting}
\FullSourceStop

At this stage everything is new, so the whole file is highlighted. From the
next step onward, only the deltas stay green.
